% Fragmento integrable. Preambulo anfitrion: amsmath, amssymb, longtable, hyperref.
\section{Compatibility of readers and recovery of the parameter}
\label{hmtfg:fragmento:compatibilidad-de-lectores-y-restitucion-del-parametro}

\subsection{Compatibility of readers, minima and recovery of history}
\label{hmtfg:escalado:13}

The prolongation operations and ordered readers preserve the information required for root selection. The following identities specify their domains and composition.

\subsubsection{Recovered forward rule}
\label{hmtfg:escalado:13.1}

The deterministic update starts from the R36 boundary state with character \(\chi\) and internal remainder. The update is
\begin{equation}
\mathcal U_K^\chi=
\operatorname{Upd}_K^\chi\circ
\operatorname{Tra}_K^\chi\circ
\operatorname{Sel}_K^\chi,
\qquad
N_{K+1}=729N_K+d_K^\chi.
\label{hmtfg:escalado:eq:26}
\end{equation}
The digit and the update depend on the preceding state. Two forward sections with the same complete R36 state coincide by induction. The subrelation \(\operatorname{Ext}^{\chi,\rightarrow}\) is the graph of that update; the total relation \(\operatorname{Ext}\) admits multiple continuations depending on the data and characters.

Compatible prolongations transport the regional readers and the joint signature over the same prefixes.

These rules transport orientation, remainder, quotient, carry, boundary and memory. Their application to the cascade must also preserve ordered selection of the critical reader’s parameter.

\subsubsection{Independence of the antecedent constants from the free parameter}
\label{hmtfg:escalado:13.2}

Let \(B\) be the HMT state providing the wheel and its antecedent constants. Write \(\mathcal C(B)\) for the set of those readouts and \(u_B\) for the critical profile. In the domain of the critical reader, the subsequent construction is
\begin{equation}
(B,\lambda)\longmapsto f_{B,\lambda}(x)=1-\lambda u_B(x),
\qquad 0<\lambda\le 2/u_B(1).
\label{hmtfg:escalado:eq:27}
\end{equation}
For the same \(B\), the reading of the preceding constants factors through projection onto the first component:
\begin{equation}
\widetilde{\mathcal C}(B,\lambda)
=\mathcal C\circ\operatorname{pr}_B(B,\lambda)
=\mathcal C(B).
\label{hmtfg:escalado:eq:28}
\end{equation}
Consequently, two distinct values of \(\lambda\) preserve the same antecedent readouts. The equality is a property of this causal construction: the profile is fixed before \(\lambda\) is varied, and the factor \(\pi_{\mathrm{HMT}}\) cancels in its quadratic normalization.

\textbf{Lemma.} Suppose a condition \(A\) is evaluated solely on those antecedent readouts. For every fixed \(B\),
\begin{equation}
\{\lambda\in\Lambda:A(\widetilde{\mathcal C}(B,\lambda))\}
=
\begin{cases}
\Lambda,&A(\mathcal C(B))\text{ holds},\\
\varnothing,&A(\mathcal C(B))\text{ fails}.
\end{cases}
\label{hmtfg:escalado:eq:29}
\end{equation}
\textbf{Proof.} Substituting \eqref{hmtfg:escalado:eq:28} makes the predicate independent of \(\lambda\). Therefore all parameters in the domain satisfy it, or none do. This includes comparison of the antecedent constants with metrological intervals.

Uniqueness of the section \(B\) and the subsequent restriction \(f_{B,\lambda}^{2^n}(0)=0\) produce a solution fiber in \(\lambda\). The first-root rule acts on this fiber through the return and its itinerary. Selection and its prolongation are proved in the following results.

This result maintains the HMT order: metrology compares previously constructed outputs, and root selection is proved within the reader that produces it.

\subsubsection{Lemma on ordered transport of the selector}
\label{hmtfg:escalado:13.3}

Let \(P_n\) be an ordered set of genealogical candidates with minimum \(p_n^*\). Let \(Z_n\) be the corpus’s original set of candidate roots, with the real order, and let \(L_n:P_n\to Z_n\) be a monotone reader. Its image is said to be \textbf{coinitial} when
\begin{equation}
\forall z\in Z_n\quad\exists p\in P_n:
L_n(p)\le z.
\label{hmtfg:escalado:eq:30}
\end{equation}
Then
\begin{equation}
\boxed{L_n(p_n^*)=\min Z_n
\quad\Longleftrightarrow\quad
L_n(P_n)\text{ is coinitial in }Z_n.}
\label{hmtfg:escalado:eq:31}
\end{equation}
\textbf{Proof.} If the image of the minimum is the global minimum, take \(p=p_n^*\) in \eqref{hmtfg:escalado:eq:30}. Conversely, given \(z\in Z_n\), coinitiality provides \(p\) with \(L_n(p)\le z\). Monotonicity gives \(L_n(p_n^*)\le L_n(p)\le z\). Since \(L_n(p_n^*)\in Z_n\), it is the global minimum. Surjectivity of \(L_n\) is a sufficient condition, stronger than \eqref{hmtfg:escalado:eq:30}.

The ordered isomorphism of the continuum chart satisfies reader surjectivity. The concrete application to the return and its complete set of zeros is proved in the ordered-transport section; classification and isolation then determine its dynamical continuation.

\subsubsection{Minimal selection through survival at arbitrary depth}
\label{hmtfg:escalado:13.6}

Composition with survival also produces a precise global selector. The uniform family \(\eta=0\), constructed in section~\ref{hmtfg:escalado:1}, is retained, and the theorems on realization of the infinite itinerary and preservation of its signs are used. Let
\[
I=\left[\frac1{2u_0(1)},\frac2{u_0(1)}\right],\qquad
c_j(\lambda)=f_\lambda^j(0),\qquad
\tau_j=(-1)^{v_2(j)}.
\]
The set
\[
\Lambda_\tau=\{\lambda\in I:\operatorname{sign}c_j(\lambda)=\tau_j
\text{ for every }j\ge1\}
\]
is compact and nonempty. The derivative control in that source provides
\begin{equation}
B_j=\frac{16^j(16^j-1)}{15},\qquad
|c_j(\lambda)|>\frac2{B_j}\quad(\lambda\in\Lambda_\tau).
\label{hmtfg:escalado:eq:32}
\end{equation}
Define
\begin{equation}
\varepsilon_j=\frac1{B_j},\qquad
A_N=\{\lambda\in I:\tau_jc_j(\lambda)\ge\varepsilon_j,\
1\le j\le N\}.
\label{hmtfg:escalado:eq:33}
\end{equation}
Each \(A_N\) is compact by continuity of the iterates and nonempty because it contains \(\Lambda_\tau\). Moreover,
\begin{equation}
A_{N+1}\subseteq A_N,\qquad
\bigcap_{N\ge1}A_N=\Lambda_\tau.
\label{hmtfg:escalado:eq:34}
\end{equation}
The right-to-left inclusion follows from \eqref{hmtfg:escalado:eq:32}. The reverse inclusion follows from the strictly positive margins in \eqref{hmtfg:escalado:eq:33}, which fix each sign.

\textbf{Surviving-minimum theorem.} The minima \(a_N=\min A_N\) satisfy
\begin{equation}
\boxed{a_N\uparrow a_*=\min\Lambda_\tau.}
\label{hmtfg:escalado:eq:35}
\end{equation}
\textbf{Proof.} The inclusion \eqref{hmtfg:escalado:eq:34} makes the sequence of minima increasing and bounded within \(I\); it therefore has a limit \(a\). For each \(r\), all terms \(a_N\), \(N\ge r\), belong to the closed set \(A_r\); thus \(a\in A_r\). By \eqref{hmtfg:escalado:eq:34}, \(a\in\Lambda_\tau\). For any \(\lambda\in\Lambda_\tau\), we have \(a_N\le\lambda\) at every level, and on passing to the limit, \(a\le\lambda\). This proves \eqref{hmtfg:escalado:eq:35}.

In the HMT cylindrical reader, the object \(a_*\) has compatible prefixes at every depth, with orientation, remainder and memory preserved. Its choice uses the order of the parameter reading and complete survival. Compactness proves the existence and convergence of these minima.

The parameters \(a_N\) solve sign constraints with a margin. Their boundaries may satisfy \(\tau_jc_j=\varepsilon_j\), whereas the superstable roots satisfy \(c_{2^n}=0\). Therefore \eqref{hmtfg:escalado:eq:35} establishes a global selector for infinite realization and preserves, as a distinct object, the corpus’s first-root selector.

\subsubsection{Recovery of the parameter from the complete critical history}
\label{hmtfg:escalado:13.7}

The second term of the critical history provides an exact inverse reader:
\begin{equation}
c_1(\lambda)=1,\qquad
c_2(\lambda)=1-\lambda u_0(1),\qquad
\boxed{\lambda=\frac{1-c_2(\lambda)}{u_0(1)}.}
\label{hmtfg:escalado:eq:36}
\end{equation}
The positivity \(u_0(1)>0\), proved for the profile, makes division valid. Consequently, two identical complete critical histories have the same parameter. The sign word \(\tau\) is a reduced readout of that history; it preserves its combinatorics and omits the magnitude of \(c_2\).

Formulas \eqref{hmtfg:escalado:eq:35}–\eqref{hmtfg:escalado:eq:36} unite survival-based selection and recovery of the parameter from the critical history. Equality of its limit with the stable crossing is proved by first exit, and identification of the centers by analytic transport of their hybrid classes.
